\documentclass[aspectratio=169,10pt]{beamer}

\usepackage[utf8]{inputenc}
\usepackage{amsmath,amssymb,amsfonts}
\usepackage{mathtools}
\usepackage{xcolor}
\usepackage{listings}
\usepackage{graphicx}
\usepackage{hyperref}
\usepackage{tikz}

\usetheme{Madrid}
\usecolortheme{seahorse}

\definecolor{codegreen}{rgb}{0, 0.6, 0}
\definecolor{codegray}{rgb}{0.5, 0.5, 0.5}
\definecolor{codepurple}{rgb}{0.58, 0, 0.82}
\definecolor{backcolour}{rgb}{0.95, 0.95, 0.92}

\lstdefinestyle{pythonstyle}{
    backgroundcolor=\color{backcolour},
    commentstyle=\color{codegreen}\itshape,
    keywordstyle=\color{codepurple}\bfseries,
    numberstyle=\tiny\color{codegray},
    stringstyle=\color{codepurple},
    basicstyle=\ttfamily\scriptsize,
    breakatwhitespace=false,
    breaklines=true,
    captionpos=b,
    keepspaces=true,
    numbers=left,
    numbersep=5pt,
    showspaces=false,
    showstringspaces=false,
    showtabs=false,
    tabsize=4,
    frame=single
}

\lstset{style=pythonstyle}

\graphicspath{{figures/}}
\setbeamertemplate{footline}[frame number]

\title{Continuity \& Measurability of Multifunctions}
\subtitle{Pathak Chapter 10b: Applications to Integral Inclusions}
\author{H. K. Pathak}
\institute{Nonlinear Analysis and Fixed Point Theory}
\date{2018}

\begin{document}

\frame{\titlepage}

\begin{frame}{Chapter Overview}
    \begin{itemize}
        \item \textbf{Topic}: Fixed point theorems for multifunctions and their applications
        \item \textbf{Focus}: Continuous and measurable multifunctions
        \item \textbf{Applications}: Nonlinear integral inclusions and differential inclusions
        \vspace{0.3cm}
        \item \textbf{Key Concepts}:
        \begin{enumerate}
            \item Continuity properties of set-valued maps
            \item Measurability of multifunctions
            \item Fixed point theorems
            \item Integral inclusions
            \item Hammerstein type equations
        \end{enumerate}
    \end{itemize}
\end{frame}

\begin{frame}[allowframebreaks]{1. Introduction to Multifunctions}
    \textbf{Definition}: A multifunction (set-valued map) $F: X \to 2^Y$ assigns to each $x \in X$ a set $F(x) \subseteq Y$.

    \vspace{0.4cm}
    \textbf{Key Distinctions}:
    \begin{itemize}
        \item Single-valued functions map points to points
        \item Multifunctions map points to sets
        \item Can express inclusion problems elegantly
    \end{itemize}

    \vspace{0.4cm}
    \textbf{Examples of Multifunctions}:
    \begin{enumerate}
        \item Normal cone to a convex set
        \item Subdifferential of a convex function: $\partial f(x) = \{y \in X^* : \langle y, z-x \rangle \leq f(z) - f(x) \text{ for all } z\}$
        \item Solutions set of differential inclusions
    \end{enumerate}

    \vspace{0.4cm}
    \textbf{Motivation}:
    \begin{itemize}
        \item Model uncertainty and set-valued constraints
        \item Study existence of solutions to inclusions rather than equations
        \item Unified framework for many problems
    \end{itemize}
\end{frame}

\begin{frame}[allowframebreaks]{2. Continuity of Multifunctions}
    \textbf{Definition (Upper Semicontinuity)}:
    A multifunction $F: X \to 2^Y$ is upper semicontinuous (u.s.c.) at $x_0$ if for every open set $U$ containing $F(x_0)$, there exists a neighborhood $V$ of $x_0$ such that $F(x) \subseteq U$ for all $x \in V$.

    \vspace{0.4cm}
    \textbf{Definition (Lower Semicontinuity)}:
    $F$ is lower semicontinuous (l.s.c.) at $x_0$ if for every open set $U$ intersecting $F(x_0)$, there exists a neighborhood $V$ of $x_0$ such that $F(x) \cap U \neq \emptyset$ for all $x \in V$.

    \vspace{0.4cm}
    \textbf{Complete Continuity}:
    $F$ is completely continuous if it maps bounded sets to sets with compact closure.

    \vspace{0.4cm}
    \textbf{Key Properties}:
    \begin{itemize}
        \item \textbf{Closed-valuedness}: $F$ is u.s.c. and locally bounded $\Rightarrow$ $F$ is closed-valued
        \item \textbf{Compactness}: If $X$ is compact and $F$ is u.s.c. with compact values $\Rightarrow$ graph of $F$ is compact
        \item \textbf{Continuity}: $F$ is continuous iff both u.s.c. and l.s.c.
    \end{itemize}
\end{frame}

\begin{frame}[allowframebreaks]{3. Measurability of Multifunctions}
    \textbf{Definition}: A multifunction $F: \Omega \to 2^X$ (where $\Omega$ is a measure space) is measurable if:

    $$F^{-}(B) := \{s \in \Omega : F(s) \cap B \neq \emptyset\} \text{ is measurable for all open } B \subseteq X$$

    \vspace{0.4cm}
    \textbf{Equivalently}:

    $$F^{+}(K) := \{s \in \Omega : F(s) \subseteq K\} \text{ is measurable for all compact } K \subseteq X$$

    \vspace{0.4cm}
    \textbf{Measurable Selection Theorem}:
    If $F: \Omega \to 2^X$ is measurable with nonempty closed values, then there exists a measurable selection $f: \Omega \to X$ such that $f(s) \in F(s)$ for almost all $s \in \Omega$.

    \vspace{0.4cm}
    \textbf{Practical Importance}:
    \begin{itemize}
        \item Allows integration of set-valued maps via selections
        \item Essential for variational methods
        \item Links multifunctions to classical measure theory
    \end{itemize}
\end{frame}

\begin{frame}[allowframebreaks]{4. Fixed Point Theorems for Multifunctions}
    \textbf{Theorem 5.167 (Kakutani-Fan-Glicksberg)}:
    Let $K$ be a nonempty compact convex subset of a Banach space $X$. If $F: K \to 2^K$ is u.s.c. with nonempty compact convex values, then $F$ has a fixed point, i.e., there exists $x \in K$ such that $x \in F(x)$.

    \vspace{0.4cm}
    \textbf{Important Consequence}:
    For applications to integral inclusions, we can construct operator $A: C \to 2^C$ (or $A: C \to K(C)$) such that fixed points correspond to solutions.

    \vspace{0.4cm}
    \textbf{Completeness Conditions}:
    \begin{itemize}
        \item $F$ is upper semicontinuous
        \item $F$ maps into nonempty compact convex sets
        \item Domain is nonempty compact convex subset of Banach space
    \end{itemize}

    \vspace{0.4cm}
    \textbf{Application Strategy}:
    \begin{enumerate}
        \item Formulate inclusion as multifunction equation
        \item Verify continuity and convexity conditions
        \item Apply fixed point theorem
        \item Construct solution from fixed point
    \end{enumerate}
\end{frame}

\begin{frame}[allowframebreaks]{5. Nonlinear Integral Inclusions}
    \textbf{General Form}:
    $$x(t) \in \lambda(t) + \int_0^t [a(s,x(s)) + f(s, x(s), u(s))] ds$$
    for $t \in [0,T]$, where:
    \begin{itemize}
        \item $\lambda \in C[0,T]$ is given
        \item $a: [0,T] \times \mathbb{R} \to \mathbb{R}$
        \item $f: [0,T] \times \mathbb{R}^2 \to 2^{\mathbb{R}}$ is a multifunction
        \item $u: [0,T] \to \mathbb{R}$ is a measurable function
    \end{itemize}

    \vspace{0.4cm}
    \textbf{Solution Space}:
    $E = C([0,T], \mathbb{R})$ with the supremum norm $\|x\|_\infty = \sup_{t \in [0,T]} |x(t)|$.

    \vspace{0.4cm}
    \textbf{Operator Formulation}:
    Let $A: C \to 2^C$ be defined by:
    $$Ax = \left\{ y \in C[0,T] : y(t) \in \lambda(t) + \int_0^t [a(s,x(s)) + f(s,x(s),u(s))] ds \right\}$$

    \vspace{0.4cm}
    \textbf{Solution}: $x$ is a solution if $x \in Ax$ (fixed point of $A$).
\end{frame}

\begin{frame}[allowframebreaks]{6. Continuous Solutions (Section 10.1)}
    \textbf{Problem 10.1}: Find continuous solutions to nonlinear integral inclusions

    \vspace{0.3cm}
    \textbf{Theorem (Continuous Solution Existence)}:
    Consider the inclusion
    $$x(t) \in \lambda(t) + \int_0^t [a(s,x(s)) + f(s,x(s),u(s))] ds, \quad t \in [0,T]$$

    \textbf{Hypotheses}:
    \begin{enumerate}
        \item $\lambda \in C[0,T]$, $a, f$ satisfy Carathéodory conditions
        \item Appropriate growth conditions: $|a(s,x)| \leq k_1(s)$, $|f(s,x)| \leq k_2(s)$
        \item Measurable selections exist for $f$
        \item Upper semicontinuity conditions on the operator
    \end{enumerate}

    \vspace{0.4cm}
    \textbf{Key Steps}:
    \begin{itemize}
        \item Construct solution set in bounded subset $C \subseteq C[0,T]$
        \item Apply Kakutani fixed point theorem
        \item Use Nemytskii operator properties
        \item Verify multifunction has fixed point
    \end{itemize}

    \vspace{0.4cm}
    \textbf{Result}: Under suitable conditions, there exists at least one continuous solution.
\end{frame}

\begin{frame}[allowframebreaks]{7. Nemytskii Operator}
    \textbf{Definition}: Let $f: \Omega \times \mathbb{R}^n \to \mathbb{R}$ be a Carathéodory function (i.e., continuous in $x$ for a.e. $s$, measurable in $s$ for all $x$). The Nemytskii operator is:

    $$N_f(u)(s) = f(s, u(s))$$

    \vspace{0.4cm}
    \textbf{Theorem 1.42 (Nemytskii Property)}:
    If $\mu(\Omega) < \infty$ and $f$ is Carathéodory, then:

    If $\{x_k(s)\}$ converges to $x_0(s)$ in measure, then $\{N_f x_k(s)\}$ converges to $N_f x_0(s)$ in measure.

    \vspace{0.4cm}
    \textbf{Continuity Theorem}:
    Suppose $N_f$ maps $L_p$ into $L_q$ (where $\tfrac{1}{p} + \tfrac{1}{q} = 1$). Then:

    $$N_f \text{ is continuous and bounded.}$$

    \vspace{0.4cm}
    \textbf{Applications}:
    \begin{itemize}
        \item Validates that solution operators are well-defined
        \item Ensures compositions of operators preserve measurability
        \item Critical for weak solutions in PDE theory
    \end{itemize}
\end{frame}

\begin{frame}[allowframebreaks]{8. Hammerstein Integral Inclusions (Section 10.2)}
    \textbf{Form}:
    $$x(t) \in \lambda(t) + \int_0^t k(t,s)g(s, x(s)) ds$$

    where $k: [0,T]^2 \to \mathbb{R}$ is a kernel, $g: [0,T] \times \mathbb{R} \to 2^{\mathbb{R}}$ is a multifunction.

    \vspace{0.4cm}
    \textbf{Key Assumptions}:
    \begin{enumerate}
        \item Kernel $k(t,s)$ satisfies uniform boundedness:
        $$\sup_{t \in [0,T]} \int_0^T |k(t,s)| ds < \infty$$
        \item Multifunction $g$ has measurable selections
        \item Growth condition: $|g(s,x)| \leq M|x| + N$ for some constants
    \end{enumerate}

    \vspace{0.4cm}
    \textbf{Existence Theorem}:
    Under the above hypotheses and contractivity conditions, there exists a continuous solution.

    \vspace{0.4cm}
    \textbf{Proof Strategy}:
    \begin{itemize}
        \item Construct operator $M_{x,\sigma}(t) := F(t, V(x_{\sigma,\lambda}(t)))$
        \item Define selection operators $T(\sigma)$
        \item Verify contractivity via Hypothesis 10.2
        \item Apply Banach fixed point theorem
    \end{itemize}
\end{frame}

\begin{frame}[allowframebreaks]{9. Filippov Type Existence Theorem (Section 10.3)}
    \textbf{Problem}: Prove existence for nonconvex integral inclusions using Filippov approach.

    \vspace{0.3cm}
    \textbf{Setup}:
    $$x(t) \in \lambda(t) + \int_0^t [a(t,s)g(t, s(u(s)) + f(t,s,u(s))] ds$$

    where $a, g, f$ are given and may be nonconvex.

    \vspace{0.3cm}
    \textbf{Key Theorem}:
    \begin{enumerate}
        \item Construct auxiliary set-valued maps $\tilde{F}$ and $\tilde{T}$
        \item Verify measurability with nonempty closed values
        \item Show $\tilde{T}(\sigma)$ is a contraction in $L_1$:
        $$H^+(T_\sigma(\sigma_1), T_\sigma(\sigma_2)) \leq \frac{1}{\sqrt{\alpha}} \|\sigma_1 - \sigma_2\|_1$$
        \item Apply Banach contraction principle
    \end{enumerate}

    \vspace{0.3cm}
    \textbf{Important Features}:
    \begin{itemize}
        \item Works for nonconvex-valued multifunctions
        \item Combines measurable selection theory with contraction principle
        \item Exponential weighting ensures contractivity
    \end{itemize}
\end{frame}

\begin{frame}[fragile]{10. Numerical Example: Integral Inclusion}
    \textbf{Problem}: Solve numerically
    $$x(t) \in 0.5 + \int_0^t 0.1 \, \sin(s) \, x(s) \, ds$$
    on $[0, 1]$ using piecewise linear approximations.

    \begin{lstlisting}[language=Python, caption=Numerical Integration]
import numpy as np
from scipy.integrate import odeint
import matplotlib.pyplot as plt

def integrand(x, t):
    """Approximate the kernel integral"""
    return 0.1 * np.sin(t) * x

def solve_inclusion():
    # Initial condition and time points
    x0 = 0.5
    t = np.linspace(0, 1, 100)

    # Solve ODE: dx/dt = 0.1*sin(t)*x
    # Solution: x(t) = x0 * exp(integral)
    solution = odeint(integrand, x0, t)

    return t, solution.flatten()

t, x = solve_inclusion()
print(f"Solution at t=1: {x[-1]:.6f}")
    \end{lstlisting}
\end{frame}

\begin{frame}[fragile]{11. Numerical Implementation (Continued)}
    \begin{lstlisting}[language=Python, caption=Verification and Visualization]
def verify_solution(t, x):
    """Verify the integral inclusion"""
    dt = t[1] - t[0]

    # Compute integral term
    integrand_vals = 0.1 * np.sin(t[:-1]) * x[:-1]
    integral = np.cumsum(integrand_vals) * dt

    # Compute inclusion RHS
    rhs = 0.5 + integral

    # Residual
    error = np.abs(x[1:] - rhs)

    return error.max(), error

# Verify and plot
error_max, error = verify_solution(t, x)
print(f"Maximum error: {error_max:.2e}")

plt.figure(figsize=(10, 6))
plt.plot(t, x, 'b-', linewidth=2, label='Solution')
plt.grid(True, alpha=0.3)
plt.xlabel('Time t')
plt.ylabel('x(t)')
plt.title('Integral Inclusion Solution')
plt.legend()
plt.tight_layout()
    \end{lstlisting}
\end{frame}

\begin{frame}[allowframebreaks]{12. Key Theorems Summary}
    \textbf{Theorem 1}: (Upper Semicontinuity and Complete Continuity)

    If $A: C \to K(C)$ is u.s.c. and completely continuous (maps bounded sets to precompact sets), then $A$ has a fixed point.

    \vspace{0.4cm}
    \textbf{Theorem 2}: (Measurable Selections)

    If $F: \Omega \to 2^X$ is measurable with nonempty closed values, there exists a measurable selection $f: \Omega \to X$ with $f(s) \in F(s)$ a.e.

    \vspace{0.4cm}
    \textbf{Theorem 3}: (Nemytskii Operator Boundedness)

    Under Carathéodory conditions with appropriate growth, the Nemytskii operator is continuous and bounded.

    \vspace{0.4cm}
    \textbf{Theorem 4}: (Contraction Principle for Multifunctions)

    If $T: L_1 \to 2^{L_1}$ is a contraction in the Hausdorff metric with compression ratio $\rho < 1$, then $T$ has a unique fixed point.
\end{frame}

\begin{frame}[allowframebreaks]{13. Key Definitions and Notation}
    \textbf{Standard Spaces}:
    \begin{itemize}
        \item $C[0,T]$: Continuous functions on $[0,T]$ with $\|x\|_\infty = \sup_{t} |x(t)|$
        \item $L_p[0,T]$: Lebesgue integrable functions with $\|x\|_p = \left(\int_0^T |x(t)|^p dt\right)^{1/p}$
        \item $2^X$: Power set of $X$ (all subsets)
        \item $K(C)$: Family of nonempty compact subsets of $C$
    \end{itemize}

    \vspace{0.3cm}
    \textbf{Multifunction Properties}:
    \begin{itemize}
        \item U.S.C. (upper semicontinuous): Inverse images of open sets are open
        \item L.S.C. (lower semicontinuous): For open $U$ intersecting $F(x)$, nearby points have image intersecting $U$
        \item Measurable: Preimage of open set is measurable
    \end{itemize}

    \vspace{0.3cm}
    \textbf{Distances}:
    \begin{itemize}
        \item Hausdorff distance: $H(A,B) = \max\{\sup_{a \in A} d(a,B), \sup_{b \in B} d(b,A)\}$
        \item Used to measure distance between sets
    \end{itemize}
\end{frame}

\begin{frame}[allowframebreaks]{14. Applications and Extensions}
    \textbf{ODE Applications}:
    \begin{itemize}
        \item Differential inclusions: $x'(t) \in F(t, x(t))$
        \item Uncertain differential equations
        \item Optimal control with set-valued constraints
    \end{itemize}

    \vspace{0.3cm}
    \textbf{PDE Applications}:
    \begin{itemize}
        \item Parabolic differential inclusions
        \item Evolution inclusions with monotone operators
        \item Systems with hysteresis effects
    \end{itemize}

    \vspace{0.3cm}
    \textbf{Advanced Topics}:
    \begin{itemize}
        \item Parametric inclusions and sensitivity analysis
        \item Viability theory (solutions staying in given set)
        \item Multivalued stochastic differential equations
    \end{itemize}

    \vspace{0.3cm}
    \textbf{Practical Uses}:
    \begin{itemize}
        \item Control theory under uncertainty
        \item Variational inequalities
        \item Image processing and optimization
    \end{itemize}
\end{frame}

\begin{frame}[allowframebreaks]{15. Summary and Conclusions}
    \textbf{Main Concepts Covered}:

    \begin{enumerate}
        \item \textbf{Multifunctions}: Set-valued maps generalize single-valued functions
        \item \textbf{Continuity}: Upper and lower semicontinuity are key topological properties
        \item \textbf{Measurability}: Essential for integrating multifunctions
        \item \textbf{Fixed Points}: Kakutani-Fan-Glicksberg theorem applies to u.s.c. maps
        \item \textbf{Integral Inclusions}: Applications to nonlinear equations with uncertainty
    \end{enumerate}

    \vspace{0.4cm}
    \textbf{Key Theoretical Results}:
    \begin{itemize}
        \item Existence of continuous solutions to integral inclusions
        \item Hammerstein type existence theorems
        \item Filippov's approach for nonconvex multifunctions
        \item Integration of multifunctions via selections
    \end{itemize}

    \vspace{0.4cm}
    \textbf{Why This Matters}:
    \begin{itemize}
        \item Provides rigorous framework for set-valued problems
        \item Connects topology, measure theory, and functional analysis
        \item Essential for modern applications in control and optimization
    \end{itemize}
\end{frame}

\begin{frame}{References and Further Study}
    \textbf{Key References}:
    \begin{enumerate}
        \item Pathak, H.K. (2018). \textit{An Introduction to Nonlinear Analysis and Fixed Point Theory}. Springer.
        \item Aubin, J.P. and Frankowska, H. (1990). \textit{Set-Valued Analysis}. Birkhauser.
        \item Deimling, K. (1992). \textit{Multivalued Differential Equations}. De Gruyter.
        \item Castaing, C. and Valadier, M. (1977). \textit{Convex Analysis and Measurable Multifunctions}. Springer.
    \end{enumerate}

    \vspace{0.4cm}
    \textbf{Topics for Further Exploration}:
    \begin{itemize}
        \item Sensitivity analysis for parametric inclusions
        \item Fuzzy differential equations
        \item Stochastic multivalued equations
        \item Applications to game theory and economics
    \end{itemize}
\end{frame}

\end{document}
